\subsection{LATTICES}

DEFINITION 8. An ordered set E is said to be a lattice if every subset consisting of two elements of E has a least upper bound and a greatest lower bound in E.

Every product of lattices is a lattice; this follows from the condition for the existence of a least upper bound in a product of ordered sets (no. 9, Proposition 8). The set of subsets of a set A, ordered by inclusion, is a lattice because the union and intersection of two subsets of A are again subsets of A.

\minorheading{Examples}

* (1) The set of integers \(\geqslant 1\) , ordered by the relation `` \(m\) divides \(n\) '' between \(m\) and \(n\) , is a lattice; the least upper bound of \(\{m, n\}\) is the l.c.m. of \(m\) and \(n\) , and the greatest lower bound is their h.c.f.

\begin{enumerate}
\def\labelenumi{(\arabic{enumi})}
\setcounter{enumi}{1}
\item
  The set of subgroups of a group G, ordered by inclusion, is a lattice.
\item
  The set of topologies on a set A, ordered by the relation `` \(\mathfrak{G}\) is coarser than \(\mathfrak{G}'\)'' between \(\mathfrak{G}\) and \(\mathfrak{G}'\) , is a lattice. (General Topology, Chapter I, § 2).
\item
  The set \(\mathcal{F}(\mathbf{I},\mathbf{R})\) of all real-valued functions defined on an interval \(\mathbf{I}\) of \(\mathbf{R}\) is a lattice with respect to the order relation \(f\leqslant g\) (no. 4), and as such is isomorphic to the product \(\mathbf{R}^{\mathrm{I}}\) . *
\end{enumerate}

Remark. A lattice is obviously both left and right directed. But an ordered set which is both left and right directed is not necessarily a lattice. * An example of the latter is the set of mappings \(x \to p(x)\) of \(\mathbf{R}\) into itself, where \(p\) is a polynomial in \(\mathbf{R}[\mathbf{X}]\) , this set being ordered by the relation \(p \leqslant q\) (no. 4). *

